% ============================== CHAPTER 9 ===================================
\chapter{Conclusion, Applications and Future Work}
\label{ch:conclusion}

\section{Conclusion}
\label{sec:conclusion}

\setu{} set out to put high-quality Indic machine translation on a device that
has no network connection, and it does. Twenty-two translation directions
covering all twenty-two scheduled languages run from a 103.9\,MB artifact each, complete
a sentence in 105 to 202 milliseconds of commodity CPU time, and produce correct,
idiomatic output with every network socket closed. Fourteen of those directions
reach at least 80\,\% of their teacher's BLEU; the remaining eight fall between
0.755 and 0.787, and none is broken.

\subsection{Outcome Against Each Objective}
\label{subsec:outcomes}

\begin{bullets}
\item \textbf{Ratio $\geq 0.80$.} 14 of 22 directions; mean 0.838, median 0.855, best 0.955. \emph{Verdict: partially met.}
\item \textbf{Latency $<500$\,ms.} Mean 105--202\,ms; only one p90 outlier. \emph{Verdict: met.}
\item \textbf{Size $\leq 200$\,MB.} 103.9\,MB for every direction, 48\,\% under budget. \emph{Verdict: met.}
\item \textbf{Fully offline.} Verified with all sockets disabled, for every direction. \emph{Verdict: met.}
\item \textbf{Test the distillation hypothesis.} Hypothesis refuted under controlled conditions: \seqkd{} 17.09 BLEU against 11.10 for DPO. \emph{Verdict: met.}
\item \textbf{Coverage and extensibility.} All 22 languages both ways; 44 models serve 506 directions via English pivot; adding a language is a configuration change. \emph{Verdict: met.}
\end{bullets}

\subsection{Principal Contributions}
\label{subsec:contributions}

\textbf{An empirical finding, obtained by refuting the project's own
hypothesis.} For compact on-device Indic machine translation, sequence-level
knowledge distillation from teacher output decisively outperforms
preference-based distillation on human references: 17.09 against 11.10 BLEU at
matched student size on identical data, an improvement of 54\,\%. The
explanation is that reference noise rather than the training objective is the
binding constraint. The evidence for that explanation is that reference-trained
students plateau near a ratio of 0.66 regardless of data volume, whereas
distilled students continue improving through 0.607, 0.764 and 0.806 as data
scales. The result is strengthened, not weakened, by the fact that the DPO
implementation works correctly: its margin accuracy reaches 1.0 and it improves
chrF exactly as its chrF-based ranking predicts. The preference signal is
learned faithfully; it is simply not the signal a 52-million-parameter student
most needs.

\textbf{A working offline system, not a prototype.} Twenty-two quantised models
are published on a public release, fetchable without a login, and served through
a REST API, an offline progressive web application, a browser client and a
command-line tool that all share one routing layer. A fresh machine can be
running translations in five commands.

\textbf{Quadratic coverage from linear storage.} The English-pivot router turns
44 trained models into 506 usable directions at a measured two-hop cost of about
50 to 65 milliseconds. At 104\,MB per model, training direct models for every
pair would have required more than 11\,GB of device storage for the same
coverage.

\textbf{A reproducible engineering record.} The project documents three distinct
training failure modes (empty output from a label-encoding mismatch,
source-agnostic output from unmasked padding in the loss, and mid-training
collapse from missing gradient clipping), each with its diagnosis and fix. It
further documents two learning-rate collapses and their complete recovery, and
establishes that supervised training loss (1.5 to 2.3 healthy, 3 to 4 collapsed)
is a cheap early indicator available long before evaluation. These failures are
routinely encountered and rarely published.

\subsection{Limitations}
\label{sec:limits}

The project states its limitations plainly, because each of them bounds how
the results should be read.

\begin{bullets}
\item \textbf{Eight directions fall short of the quality target,} at ratios
between 0.755 and 0.787. Their supervised losses cluster at 2.00 to 2.32, the
upper end of the healthy band, which indicates under-convergence rather than
instability.

\item \textbf{The teacher caps the ceiling.} The ungated distilled-200M teacher
scores only 27 to 29 BLEU on these slices, so a student at 0.806 of that sits
at roughly four fifths of a model that is itself well below the gated
one-billion-parameter flagship. Raising the ceiling requires teacher access,
not more student training.

\item \textbf{The reported ratios are not benchmark numbers.} They come from an
internal 500-sentence Samanantar development slice, and the teacher's own score
varies between slices (27.08 to 28.91 for Hindi\arrowto{}English alone). BLEU
also understates the agglutinative languages: English\arrowto{}Malayalam passes
at a ratio of 0.863 on 5.2 BLEU. The figures are internally comparable but
should not be set beside published IndicTrans2 results without FLORES-200
rescoring under COMET.

\item \textbf{ARM latency is projected, not measured.} All timings were taken
on x86 CPUs. The target is ARM Cortex-A55 class hardware, and although the
margin is large, roughly 150\,ms against a 500\,ms budget, the figure remains
unverified on the intended silicon.
\end{bullets}

\section{Applications}
\label{sec:applications}

\textbf{Rural healthcare.} A doctor or ASHA worker can converse with a patient
across a language boundary on a handset with no connectivity. The privacy
property matters as much as the offline one: a patient's described symptoms
never leave the device.

\textbf{Defence and border operations.} Patrol units operating beyond cellular
coverage can translate locally. Because there is no transmission, there is no
signal to intercept and no dependency on infrastructure that may be unavailable
or compromised.

\textbf{Education.} Multilingual tutoring tools can serve students in their own
language at zero marginal cost per query, which is what makes deployment across
many schools economically feasible.

\textbf{Government service delivery.} Forms, notices and entitlement documents
can be translated at a gram panchayat counter without the counter needing a
working internet connection or a commercial API subscription.

\textbf{Disaster relief.} Field workers can communicate with affected
communities precisely during the window when networks are down, the moment
when a cloud translator is guaranteed to be unavailable.

\textbf{Privacy-sensitive documents.} Legal, medical and financial text can be
translated without disclosure to a third party, which is a structural guarantee
here rather than a contractual one.

\textbf{As a research artifact.} The distillation pipeline is not
Indic-specific. Any language family with a strong teacher and a noisy parallel
corpus can use the same configuration-driven process, and the reported
comparison gives practitioners a concrete data point on which training signal to
choose.

\section{Future Scope}
\label{sec:future}

The immediate work is to rescore every direction on FLORES-200 with COMET,
which converts the internal ratios into publishable, comparable numbers and may
reclassify several agglutinative-language near-misses as passes; the
infrastructure already exists in \texttt{setu-eval}. Alongside it, the eight
near-misses should be closed with the three levers their under-convergence
signature points to, in increasing order of cost: beam search at evaluation
time, additional supervised epochs, and more distilled data. Latency should
then be measured on real ARM hardware rather than projected from x86, and the
cost of the second hop in pivoted Indic-to-Indic translation quantified against
real references. Beyond that, the single most effective remaining lever on
quality is a stronger teacher: at 0.806 against a 27-BLEU teacher the student
is near its current ceiling, and the gated one-billion-parameter IndicTrans2
flagship would lift the entire curve. The All-in-One KD and Elastic Weight
Consolidation mechanisms are implemented and tested but have not been exercised
on full multilingual data; consolidating the per-direction students into one
multilingual model would cut device storage from $n \times 104$\,MB to a single
artifact, at the risk of the cross-lingual interference that per-direction
training currently avoids. A runtime with genuine 4-bit kernels would halve the
artifact again, and the architecture is already pair-agnostic, so widening
coverage beyond the Eighth Schedule to the major non-scheduled languages and
regional dialects remains a configuration and data change rather than a
rewrite.

\section{Closing Remarks}
\label{sec:closing}

The project's most valuable outcome was not the one it set out to obtain. \setu{}
was proposed as a demonstration that preference distillation would produce better
compact translators, and the controlled experiment showed the opposite. The
correct response to that was neither to discard the DPO work nor to present the
system as though the hypothesis had held, but to report what was measured,
rebuild the system around the method that won, and treat the refutation as the
contribution.

What remains is a translator that fits in 104 megabytes, answers in a sixth of a
second, works with the network switched off, covers eleven of India's scheduled
languages in every direction between them, and never sends a user's words
anywhere. The eight directions that fall short of target are documented as
falling short. That, rather than a uniformly successful scorecard, is what the
work actually produced.
